% Source PDF pages 85--87; manual pages 2-56--2-58.
\subsection{Aerodynamic Forces}\label{sec:aerodynamic-forces}

Aerodynamic force vectors are calculated from aerodynamic coefficients supplied in
a table file or problem file, as described in
\cref{sec:table-types,sec:block-aero}. A coefficient is converted to force by
multiplying it by dynamic pressure $q$ and aerodynamic reference area
$S_{\mathrm{ref}}$.

TAOS accepts three coefficient systems: axial and normal coefficients
$(C_A,C_N)$; lift, drag, and side-force coefficients $(C_L,C_D,C_S)$; and body
$x$-, $y$-, and $z$-axis coefficients $(C_X,C_Y,C_Z)$. Each set defines an
aerodynamic force vector. If more than one set is supplied, their force vectors are
summed to form $\sum\vec F_{\mathrm{aero}}$.

Aerodynamic coefficients are often functions of Mach number and Reynolds number.
TAOS computes those quantities together with dynamic pressure as
\begin{center}
\begin{tabularx}{0.90\linewidth}{@{}>{\ttfamily}l>{$}l<{$}X@{}}
  dynprs & q=0.7pM^2 & dynamic pressure.\\
  mach   & M=\lVert\vec V_{w\oplus}\rVert/c & Mach number.\\
  reypft & R_N=\lVert\vec V_{w\oplus}\rVert/\nu & Reynolds number per unit length.\\
\end{tabularx}
\end{center}
These variables depend on atmospheric pressure $p$, speed of sound $c$, and
kinematic viscosity $\nu$ from \cref{sec:atmosphere-model}.

\taossubhead{Axial and Normal Force Coefficients}

Axial and normal coefficients are commonly used for axisymmetric vehicles such as
reentry bodies and missiles. The axial force acts opposite the body $x$ axis, and the
normal force acts opposite the windward meridian $\phi_w$.

A body-frame unit vector in the windward-meridian direction is obtained by
projecting the wind $x$ axis onto the plane spanned by $\hat y_b$ and $\hat z_b$,
as illustrated by \cref{fig:alpha-betae-definitions}:
\begin{equation}
  \hat\phi_b
  =\frac{
      (\hat x_w\mathbin{\cdot}\hat y_b)\hat y_b
      +(\hat x_w\mathbin{\cdot}\hat z_b)\hat z_b}
    {\sqrt{
      (\hat x_w\mathbin{\cdot}\hat y_b)^2
      +(\hat x_w\mathbin{\cdot}\hat z_b)^2}}.
  \label{eq:windward-meridian-unit-vector}
\end{equation}
The axial and normal coefficients then produce
\begin{equation}
  \vec F_{\mathrm{aero}}
  =-\left(C_A\hat x_b+C_N\hat\phi_b\right)qS_{\mathrm{ref}}.
  \label{eq:aero-force-axial-normal}
\end{equation}
The first term is the axial force and the second is the normal force.
\Cref{eq:windward-meridian-unit-vector} is singular when $\hat x_w$ is aligned
with $\hat x_b$, corresponding to zero total angle of attack. TAOS sets the normal
force to zero in that case.

\taossubhead{Lift, Drag, and Side-Force Coefficients}

Lift, drag, and side-force coefficients are commonly used for aircraft. Drag acts
opposite $\hat x_w$, lift acts opposite $\hat z_w$, and side force acts in the
positive $\hat y_w$ direction. The aerodynamic force vector is therefore
\begin{equation}
  \vec F_{\mathrm{aero}}
  =-\left(C_D\hat x_w-C_S\hat y_w+C_L\hat z_w\right)
    qS_{\mathrm{ref}}.
  \label{eq:aero-force-lift-drag-side}
\end{equation}

\taossubhead{Body-Axes Force Coefficients}

Coefficients $C_X$, $C_Y$, and $C_Z$, which may be obtained directly from wind
tunnel testing, define forces aligned with the body-fixed coordinate system:
\begin{equation}
  \vec F_{\mathrm{aero}}
  =\left(C_X\hat x_b+C_Y\hat y_b+C_Z\hat z_b\right)
    qS_{\mathrm{ref}}.
  \label{eq:aero-force-body-axes}
\end{equation}

\taossubhead{Aerodynamic-Force Evaluation}

The function \texttt{force\_aero} evaluates aerodynamic forces. It first converts
each input coefficient system to lift, drag, and side-force components using the
preceding definitions and coordinate transformations. If a segment contains
multiple \problemblock{*aero} data blocks, the force from each block is added to
form the total lift, drag, and side force. The final force is transformed from the wind
system to ECFC coordinates for use in the equations of motion.

The aerodynamic-coefficient output variables are
\begin{center}
\begin{tabularx}{0.86\linewidth}{@{}>{\ttfamily}l>{$}l<{$}X@{}}
  ca & C_A & axial aerodynamic force coefficient.\\
  cn & C_N & normal aerodynamic force coefficient.\\
  \addlinespace[0.25em]
  cl & C_L & total aerodynamic lift coefficient.\\
  cd & C_D & total aerodynamic drag coefficient.\\
  cs & C_S & total aerodynamic side-force coefficient.\\
  \addlinespace[0.25em]
  cx & C_X & body $x$-axis aerodynamic force coefficient.\\
  cy & C_Y & body $y$-axis aerodynamic force coefficient.\\
  cz & C_Z & body $z$-axis aerodynamic force coefficient.\\
\end{tabularx}
\end{center}
The lift, drag, and side-force outputs are always calculated. The other coefficient
outputs are defined only when their corresponding coefficient system is supplied.
For example, if the input tables define $C_A$ and $C_N$, the variables
\statevar{ca} and \statevar{cn} are defined while \statevar{cx}, \statevar{cy},
and \statevar{cz} remain undefined.

When multiple aerodynamic coefficient sets are supplied, only $C_L$, $C_D$, and
$C_S$ are guaranteed to represent the complete aerodynamic force. If every input
set uses the same coefficient convention, the corresponding source-system outputs
are also totaled correctly. If conventions are mixed, the remaining coefficient
outputs are either undefined or represent partial totals.
